{[}Title{]}*

{[}Your Name{]}

North Carolina Agricultural and Technical State University

A {[}thesis or dissertation{]} submitted to the graduate faculty

in partial fulfillment of the requirements for the degree of

{[}Degree Type goes here: e.g. Master of Science or Doctor of
Philosophy{]}

Program: {[}e.g. Electrical Engineering or Leadership Studies{]}

{[}Concentration: If your program has an official Concentration, you can
include it on this line{]}

Major Professor: Dr. {[}Major Professor's Name{]}

Greensboro, North Carolina

{[} Year {]}

\emph{*Please note: Wherever there are square brackets} {[} {]} \emph{on
these preliminary pages, please add the information relevant to your
thesis/dissertation and delete the square brackets (along with this
note).}

\emph{Template Updated 02-2026}

This is to certify that the {[}master's thesis or doctoral
dissertation{]} of

{[}Your name exactly as it appears on the title page{]}

has met the {[}thesis or dissertation{]} requirements of

North Carolina Agricultural and Technical State University

Greensboro, North Carolina

{[} Year {]}

Approved by:

{[}Add name \& delete brackets{]}

Major Professor

{[}Add name \& delete brackets{]}

Committee Member

{[}Add name \& delete brackets{]}

Committee Member

{[}Add name \& delete brackets{]}

Committee Member

{[}Add name \& delete brackets{]}

Graduate College Faculty Representative {[}for doctoral committees only;
delete this note{]}

{[}Add name \& delete brackets{]}

Department Chair

Dr. Clay S. Gloster, Jr.

Dean, The Graduate College

\emph{*Please note: Remove or add entries for committee members as
needed based on your committee makeup. If an individual has two roles
(e.g. if the Department Chair is also a committee member), then you can
either list the person in both places or combine the roles on one line
(e.g. Committee Member \& Department Chair). Delete this note.}

© Copyright by

{[}Your name here exactly as it appears on your title page. Please note:
this page is required{]}

{[}Year{]}

Abstract

The Abstract should be one or two pages. Your abstract begins on roman
numeral page iv.

Biographical Sketch

This page is required and should appear in paragraph form. The
Biographical Sketch must be written in the third person. Your first
sentence should begin with your full name (appearing exactly as it
appears on your Title Page, Signature Page, and Copyright Page).

Dedication

This page is optional and appears after your biographical sketch.

Acknowledgments

It is recommended that you include an Acknowledgments page. The
acknowledgments should be in paragraph form.

Table of Contents

\hypertarget{list-of-figures-x}{%
\subparagraph{\texorpdfstring{\protect\hyperlink{list-of-figures}{List
of Figures x}}{List of Figures x}}\label{list-of-figures-x}}

\hypertarget{list-of-tables-xi}{%
\subparagraph{\texorpdfstring{\protect\hyperlink{list-of-tables}{List of
Tables xi}}{List of Tables xi}}\label{list-of-tables-xi}}

\hypertarget{chapter-1-introduction-1}{%
\subparagraph{\texorpdfstring{\protect\hyperlink{chapter-1-introduction}{CHAPTER
1: Introduction
1}}{CHAPTER 1: Introduction 1}}\label{chapter-1-introduction-1}}

\protect\hyperlink{_Toc114562011}{1.1 Sample Level Two Heading and
Heading Information 1}

\protect\hyperlink{another-sample-level-two-heading}{1.2 Another Sample
Level Two Heading 1}

\protect\hyperlink{_Toc114562013}{1.2.1 Sample Level Three Heading and
Template Information 1}

\protect\hyperlink{_Toc114562014}{1.2.2 Sample Level Three Heading 2}

\protect\hyperlink{level-two-heading}{1.3 Level Two Heading 2}

\protect\hyperlink{_Toc114562016}{1.3.1 Level Three Heading 2}

\protect\hyperlink{_Toc114562017}{1.3.1.1 How to Insert a Figure. 2}

\protect\hyperlink{_Toc114562018}{1.3.1.1.1 A Note on Figure Captions
and the UPDATE FIELD Function. 3}

\hypertarget{chapter-2-literature-review-4}{%
\subparagraph{\texorpdfstring{\protect\hyperlink{chapter-2-literature-review}{CHAPTER
2: Literature Review
4}}{CHAPTER 2: Literature Review 4}}\label{chapter-2-literature-review-4}}

\protect\hyperlink{_Toc114562020}{2.1 Sample Level Two Heading and
Information About Inserting Tables 4}

\protect\hyperlink{_Toc114562021}{2.1.1 Sample Level 3 Heading 5}

\protect\hyperlink{_Toc114562022}{2.1.1.1 Sample Level Four Heading 5}

\protect\hyperlink{_Toc114562023}{2.1.1.1.1 Sample Level Five Heading.
5}

\hypertarget{chapter-3-methodology-6}{%
\subparagraph{\texorpdfstring{\protect\hyperlink{chapter-3-methodology}{CHAPTER
3: Methodology
6}}{CHAPTER 3: Methodology 6}}\label{chapter-3-methodology-6}}

\protect\hyperlink{sample-level-two-heading}{3.1 Sample Level Two
Heading 6}

\protect\hyperlink{sample-level-three-heading}{3.1.1 Sample Level Three
Heading 6}

\protect\hyperlink{_Toc114562027}{3.1.1.1 Sample Level Four Heading 6}

\protect\hyperlink{_Toc114562028}{3.1.1.1.1 Sample Level Five Heading 7}

\hypertarget{chapter-4-results-and-discussion-8}{%
\subparagraph{\texorpdfstring{\protect\hyperlink{chapter-4-results-and-discussion}{CHAPTER
4: Results and Discussion
8}}{CHAPTER 4: Results and Discussion 8}}\label{chapter-4-results-and-discussion-8}}

\protect\hyperlink{_Toc114562030}{4.1 Sample Level Two Heading 8}

\protect\hyperlink{_Toc114562031}{4.1.1 Sample Level Three Heading 8}

\protect\hyperlink{_Toc114562032}{4.1.1.1 Sample Level Four Heading. 8}

\protect\hyperlink{_Toc114562033}{4.1.1.1.1 Sample Level Five Heading.
8}

\hypertarget{chapter-5-conclusion-9}{%
\subparagraph{\texorpdfstring{\protect\hyperlink{chapter-5-conclusion}{CHAPTER
5: Conclusion
9}}{CHAPTER 5: Conclusion 9}}\label{chapter-5-conclusion-9}}

\hypertarget{references-10}{%
\subparagraph{\texorpdfstring{\protect\hyperlink{references}{References
10}}{References 10}}\label{references-10}}

\hypertarget{appendix-11}{%
\subparagraph{\texorpdfstring{\protect\hyperlink{appendix}{Appendix
11}}{Appendix 11}}\label{appendix-11}}

\hypertarget{list-of-figures}{%
\section{\texorpdfstring{\hfill\break
List of Figures}{ List of Figures}}\label{list-of-figures}}

\protect\hyperlink{_Toc114562037}{Figure 1: An Old, Rusty Ladder. 3}

\protect\hyperlink{_Toc114562038}{Figure 2: Another Image of a Ladder.
5}

\protect\hyperlink{_Toc114562039}{Figure 3: The Third Ladder. 6}

\hypertarget{list-of-tables}{%
\section{\texorpdfstring{\hfill\break
List of Tables}{ List of Tables}}\label{list-of-tables}}

\protect\hyperlink{_Toc111537318}{Table 1: Comparisons of Scores,
Numbers of Students, and Results Across Schools. 4}

\protect\hyperlink{_Toc111537319}{Table 2: Age Groups MPH Habits and
Traffic Ticket Received. 6}

\hypertarget{chapter-1-introduction}{%
\section{CHAPTER 1: Introduction}\label{chapter-1-introduction}}

Welcome to the template for theses and dissertations submitted to North
Carolina Agricultural and Technical State University. We hope this
template helps you format your document more easily. You may add and
subtract text to this template as needed, but please make sure you
understand how the headings, table of contents, and so on function
before deleting this information. This document is formatted to meet the
required standards for submission, and it also explains its own
functions here in the text. Please note that your paragraph text can use
left-aligned or justified style, but make sure to be consistent
throughout the entire document.

\protect\hypertarget{_Toc114562011}{}{}\textbf{1.1 Sample Level Two
Heading and Heading Information}

The sample headings have been provided for your benefit. You can copy
and paste them as needed. All five heading levels are provided in this
template. Many documents will need only two or three heading levels. We
recommend that students use a numbering system for headings. This is
helpful for keeping the reader (and the writer) oriented.

\hypertarget{another-sample-level-two-heading}{%
\subsection{1.2 Another Sample Level Two
Heading}\label{another-sample-level-two-heading}}

If you are numbering your sub-headings, use 1.1, 1.2, etc.to number
them. Notice that you should indent the first line of text under each
Level 2 and Level 3 heading.

\protect\hypertarget{_Toc114562013}{}{}\emph{\textbf{1.2.1 Sample Level
Three Heading and Template Information}}

This template also provides style settings for headings. To create your
own headings, first type your heading text in the desired location, and
then go to the HOME tab at the top of the page. On the right-hand side
of the HOME menu, you will see a number of style options. ``Normal''
corresponds to text typed throughout your paragraphs, ``Heading 1''
corresponds to Level One headings, ``Heading 2'' corresponds to Level
Two headings, and so on. So, type in your heading (including appropriate
numbering), highlight the heading text, and click on the appropriate
Heading Style. Now, when you use UPDATE FIELD on the Table of Contents,
your headings and page numbers will be added automatically. If you are
unfamiliar with these functions, test them out early in the
writing/formatting process so that they become familiar to you. (See
section 1.3.1.1.1 below for more information on the UPDATE FIELD
function.)

\protect\hypertarget{_Toc114562014}{}{}\emph{\textbf{1.2.2 Sample Level
Three Heading}}

Note the way section numbering proceeds in these sample subheadings.

\hypertarget{level-two-heading}{%
\subsection{\texorpdfstring{1.3 Level Two Heading
}{1.3 Level Two Heading }}\label{level-two-heading}}

\protect\hypertarget{_Toc114562016}{}{}\textbf{\emph{1.3.1 Level Three
Heading}}

Here is another example to show the numbering sequence. The next section
will cover figures and figure captions.

\protect\hypertarget{_Toc114562017}{}{}\textbf{1.3.1.1 How to Insert a
Figure.} You may click the INSERT tab at the top of this page and choose
the PICTURES option. Once the picture is in the text, be sure to center
it. All figures should be centered. Figure captions can be centered as
well. To insert a new figure caption, place your cursor one
(double-spaced) line below the figure. Choose the REFERENCES tab and
select INSERT CAPTION. Click OK. Once the caption is inserted, it is
your responsibility to format it consistently (we recommend that you use
``Figure \#:'' in bold with the caption text itself not in bold). You
can now return to your LIST OF FIGURES page, right click the text, and
select the UPDATE FIELD option. If you update the entire table, your new
captions will show up in the List of Figures.

\includegraphics[width=2.12153in,height=3.18264in]{media/image1.jpeg}

\protect\hypertarget{_Toc114562037}{}{}Figure : An Old, Rusty Ladder.

\protect\hypertarget{_Toc114562018}{}{}\emph{\textbf{1.3.1.1.1 A Note on
Figure Captions and the UPDATE FIELD Function.}} Please note: if you add
additional text or descriptions to the figure captions, everything you
type there will appear on your List of Figures page when you update the
field. It is recommended that you describe the figures and tables in
detail in the \emph{body} of your text and keep the figure caption
straightforward and to the point. The UPDATE FIELD function allows you
to automatically update your Table of Contents and Lists of
Figures/Tables. To use the UPDATE FIELD function, go to the Table of
Contents, List of Figures, or List of Tables. There are two methods: 1)
right-click the text of the Table of Contents of List of Figures/Tables
(it will become shaded) and click ``Update Field'' on the pop-up menu;
or 2) left-click the text of the Table/List (it will become shaded) and
press F9 (this is the ``refresh'' shortcut key on Windows devices). You
can now choose to ``Update page numbers only'' (this will automatically
check the document for changes to page numbers only) or ``Update entire
table'' (use this option if you have added new sections or
figures/tables). Now, all of your newly-added captions or section
headings should be pulled into the Table of Contents or List of
Figures/Tables automatically.

\hypertarget{chapter-2-literature-review}{%
\section{\texorpdfstring{\hfill\break
CHAPTER 2: Literature
Review}{ CHAPTER 2: Literature Review}}\label{chapter-2-literature-review}}

\protect\hypertarget{_Toc114562020}{}{}\textbf{2.1 Sample Level Two
Heading and Information About Inserting Tables}

Table captions are placed ABOVE your table. The table itself as well as
the table caption can be either flush against the left margin or
centered (choose one style and be consistent throughout the document).
You can select the INSERT tab at the top of the page and choose TABLE to
add a new table in your document. Then, use the INSERT CAPTION option
from the REFERENCES ribbon to insert a table number above the table.

\protect\hypertarget{_Toc111537318}{}{}Table : Comparisons of Scores,
Numbers of Students, and Results Across Schools.

\begin{longtable}[]{@{}lllll@{}}
\toprule
Schools & Average Score & Number of Students & Results & Total \\
\midrule
\endhead
MTR Academy & 7.2 & 156 & 98 & 78 \\
PWD School of Science & 8.7 & 56 & 45 & 17 \\
CFC High School & 5.2 & 356 & 52 & 36 \\
RCP Academy & 6.4 & 421 & 87 & 90 \\
\bottomrule
\end{longtable}

For Tables that need to be split across pages, you must split the table
at the appropriate place and then \textbf{manually} type in a caption on
the subsequent page (with the same Table \# and \emph{Cont.} as the
caption). See pgs. 6-7 below for an example. If you have a very wide
table, you may need to place it on a LANDSCAPE page. To add a table in
LANDSCAPE orientation, you must first create a section break by clicking
PAGE LAYOUT. Under the BREAKS options, insert a Section Break that will
start on the Next Page. This creates a new section on the next page that
can have different formatting. You can now change the page orientation
in the new section using the PAGE LAYOUT tab, and then insert your wider
table. You will need to create another section break to go back to
standard portrait pages.

\protect\hypertarget{_Toc114562021}{}{}\emph{\textbf{2.1.1 Sample Level
3 Heading}}

Please remember that you should always give context for your figures and
tables in the body of the text.

\includegraphics[width=1.64375in,height=2.06111in]{media/image2.jpeg}

\protect\hypertarget{_Toc114562038}{}{}Figure : Another Image of a
Ladder.

You should also have text after your figures and tables that explains
the item and provides additional clarifications.

\protect\hypertarget{_Toc114562022}{}{}\textbf{2.1.1.1 Sample Level Four
Heading.} Level four headings should be in bold and indented, with a
period after the heading. The paragraph text starts on the same line as
the heading. To add a level four heading using Styles, first type your
heading number, heading text, and the period after the heading. Then
highlight everything up to \textbf{(but not including)} the period and
choose the Heading 4 Style. This method helps ensure that your paragraph
text after the period isn't accidentally included in the Heading 4
Style.

\protect\hypertarget{_Toc114562023}{}{}\emph{\textbf{2.1.1.1.1 Sample
Level Five Heading.}} Level five headings are formatted much like level
four headings. The difference is that level five headings are bold AND
italicized.

\hypertarget{chapter-3-methodology}{%
\section{CHAPTER 3: Methodology}\label{chapter-3-methodology}}

\hypertarget{sample-level-two-heading}{%
\subsection{\texorpdfstring{ 3.1 Sample Level Two
Heading}{ 3.1 Sample Level Two Heading}}\label{sample-level-two-heading}}

Your text starts here on the line below the heading.

\hypertarget{sample-level-three-heading}{%
\subsubsection{\texorpdfstring{\emph{\textbf{3.1.1 Sample Level Three
Heading}}
}{3.1.1 Sample Level Three Heading }}\label{sample-level-three-heading}}

Your text starts here on the line below the heading.

\includegraphics[width=2.70417in,height=3.33056in]{media/image3.jpeg}

\protect\hypertarget{_Toc114562039}{}{}Figure : The Third Ladder.

Continue to explain the meaning and significance of your figures here in
the text, following the image.

\protect\hypertarget{_Toc114562027}{}{}\textbf{3.1.1.1 Sample Level Four
Heading}. Your paragraph text starts here.

\protect\hypertarget{_Toc111537319}{}{}Table : Age Groups MPH Habits and
Traffic Ticket Received.

\begin{longtable}[]{@{}lll@{}}
\toprule
Age Group & Average MPH & Number of Tickets \\
\midrule
\endhead
16-20 & 62 & 104 \\
20-30 & 74 & 187 \\
30-40 & 61 & 144 \\
\bottomrule
\end{longtable}

\textbf{Table 2:} Cont.

\begin{longtable}[]{@{}lll@{}}
\toprule
\endhead
40-50 & 60 & 132 \\
\bottomrule
\end{longtable}

In the event that a table does not fit on one page, you may continue it
onto the next page. You can create a break between the rows of your
table and insert the table number and the caption ``Cont.'' before
continuing the table. Be sure to type the Table \# and ``Cont.'' in
\textbf{manually} (if you were to use the INSERT CAPTION function, Word
would think you were adding an entirely new table).

\protect\hypertarget{_Toc114562028}{}{}\emph{\textbf{3.1.1.1.1 Sample
Level Five Heading}}. Your paragraph text begins here.

\hypertarget{chapter-4-results-and-discussion}{%
\section{CHAPTER 4: Results and
Discussion}\label{chapter-4-results-and-discussion}}

\protect\hypertarget{_Toc114562030}{}{}\textbf{4.1 Sample Level Two
Heading}

\protect\hypertarget{_Toc114562031}{}{}\emph{\textbf{4.1.1 Sample Level
Three Heading}}

\protect\hypertarget{_Toc114562032}{}{}\textbf{4.1.1.1 Sample Level Four
Heading.} Your text beings here.

\protect\hypertarget{_Toc114562033}{}{}\emph{\textbf{4.1.1.1.1 Sample
Level Five Heading.}} Your text begins here.

\hypertarget{chapter-5-conclusion}{%
\section{CHAPTER 5: Conclusion}\label{chapter-5-conclusion}}

Your text here.

\hypertarget{references}{%
\section{\texorpdfstring{\hfill\break
References}{ References}}\label{references}}

Entries in your References section (as well as your in-text citations)
must adhere to an accepted citation style. You may use APA
7\textsuperscript{th} edition, IEEE, MLA, or another citation style
appropriate to your field of study. It is NOT acceptable to mix citation
styles or make up your own citation style. Information on the
requirements of each style for in-text citations and final References
entries is readily available online. We also strongly recommend the use
of citation software such as RefWorks for managing your citations.
Bluford Library offers training on RefWorks and other topics related to
research.

Accumsan euismod dis natoque nisi viverra interdum tortor sodales porta
per est dolor leo malesuada ultricies hac cubilia nec posuere aliquam
quam eros mus lobortis quisque pede aliquam fames massa iaculis
facilisis nulla.

Conubia potenti class potenti eu nostra quisque id nonummy id, pulvinar
eu inceptos Pede nulla scelerisque accumsan senectus sit sem lacinia
ultricies vestibulum egestas mi accumsan congue praesent felis nunc,
adipiscing orci condimentum class dictumst egestas aliquet quisque elit.

Sem ac pede at leo, tempor scelerisque diam posuere at urna morbi
natoque luctus nullam cras ante. Nisl. Ligula mattis mauris nostra
maecenas aliquam felis aptent. Consequat aliquam velit leo.
Consectetuer.

Tellus rutrum penatibus litora malesuada faucibus nascetur interdum
consequat ipsum sociis ut potenti velit, consequat arcu per dui eleifend
dictumst montes dictumst. Volutpat cubilia quis nascetur nisl hac augue
fermentum phasellus, nonummy massa torquent sit viverra ad duis
elementum aptent egestas erat sem

\hypertarget{appendix}{%
\section{Appendix}\label{appendix}}

Various types of content can be placed in your appendices: code, sample
forms, or larger versions of tables and figures. If a figure or table is
placed in an Appendix section, it must use a figure or table caption
with a letter-number system, such as ``Figure A.1'' followed by ``Figure
A.2,'' etc. Appendix figures and tables do not necessarily need to be
included in the List of Figures or List of Tables.

If you have only one Appendix, simply label it ``Appendix.'' If you have
more than one, label them ``Appendix A,'' Appendix B,'' etc. In most
cases, it is appropriate to refer readers to the Appendix or Appendices
in the body of the document (e.g. ``A copy of Form X is provided in
Appendix A''). If you are not using an Appendix at all, you can remove
it from the document altogether.

.
